%\subsubsection*{I. Giao thoa sóng cơ}
\setcounter{paragraph}{0}
\subsubsection{Định nghĩa}
\newghichu{\textit{Hiện tượng giao thoa} là hiện tượng hai sóng kết hợp khi gặp nhau thì có những điểm ở đó chúng luôn luôn tăng cường lẫn nhau tạo thành cực đại, có những điểm ở đó chúng luôn luôn triệt tiêu nhau tạo thành cực tiểu.\\
Các cực đại và cực tiểu xen kẽ nhau tạo thành các vân giao thoa}
\subsubsection{Điều kiện giao thoa}
\Binhluan{Để các vân giao thoa ổn định trên mặt nước thì hai nguồn sóng phải là hai nguồn kết hợp có các đặc điểm:
\begin{itemize}
\item Dao động \textit{cùng phương}.
\item Cùng \textit{chu kì (hay tần số)}.
\item \textit{Hiệu số pha} không đổi theo thời gian
\end{itemize}}{
\circEX{1} Hai sóng do hai nguồn kết hợp phát ra gọi là \textit{hai sóng kết hợp}.\\
\circEX{2} Hai nguồn kết hợp cùng pha gọi là \textit{hai nguồn đồng bộ}.
}

\subsubsection{Giải thích hiện tượng giao thoa} 
\immini
{
\begin{itemize}
\item Sự giao thoa của các sóng trên thực chất chính là sự cộng gộp của các dao động. Tại mỗi điểm trong không gian, dao động của một phần tử môi trường chính là dao động tổng hợp của các dao động thành phần do các nguồn sóng truyền tới. 
\item Trong không gian các điểm có dao động được tăng cường nhau nhiều nhất khi các sóng thành phần đồng pha hoặc bị hạn chế nhau nhiều nhất khi các sóng thành phần có pha ngược nhau.
\end{itemize}
%\subsubsection*{II. Vị trí cực đại và cực tiểu} 
%\setcounter{paragraph}{0}
\subsubsection{Biên độ giao thoa sóng tại một điểm}}
{\begin{tikzpicture}[draw=Mapcolor, very thick,]
\foreach \x in {-2,-1,0,1,2}
{\node[above] at (\x,3){$\x$};}
\node[above] at (-2.7,3){$k=$};
\begin{scope}[shift={(0,0)}]
\tkzInit[ymin=-3,ymax=3,xmin=-3,xmax=3]\tkzClip
\tkzDefPoints{-2/0/s1, 2/0/s2}
\foreach \x in {0,2,4,6,8,10,12}
{\draw[dotted,thick] (s1) circle (0.3+0.5*\x);
\draw[dotted,thick] (s2) circle (0.3+0.5*\x);}
\foreach \x in {1,3,5,7,9,11}
{\draw (s1) circle (0.3+0.5*\x);
\draw (s2) circle (0.3+0.5*\x);}
\tkzDrawPoints[size=3,fill=Mapcolor](s1,s2)
\draw[dashed,blue] plot[domain=-2:2] ({1*cosh(\x)-0.75},{4.5*sinh(\x)});
\draw[dashed,blue] plot[domain=-2:2] ({-1*cosh(\x)+0.75},{4.5*sinh(\x)});
\draw[Mapcolor] plot[domain=-2:2] ({1*cosh(\x)-0.5},{3*sinh(\x)});
\draw[Mapcolor] plot[domain=-2:2] ({-1*cosh(\x)+0.5},{3*sinh(\x)});
\draw[dashed,blue] plot[domain=-2:2] ({1*cosh(\x)-0.25},{2.2*sinh(\x)});
\draw[dashed,blue] plot[domain=-2:2] ({-1*cosh(\x)+0.25},{2.2*sinh(\x)});
\draw[Mapcolor] plot[domain=-2:2] ({1*cosh(\x)},{1.75*sinh(\x)});
\draw[Mapcolor] plot[domain=-2:2] ({-1*cosh(\x)},{1.75*sinh(\x)});
\draw[dashed,blue] plot[domain=-2:2] ({1*cosh(\x)+0.25},{1.35*sinh(\x)});
\draw[dashed,blue] plot[domain=-2:2] ({-1*cosh(\x)-0.25},{1.35*sinh(\x)});
\draw[Mapcolor] plot[domain=-2:2] ({1*cosh(\x)+0.5},{1.01*sinh(\x)});
\draw[Mapcolor] plot[domain=-2:2] ({-1*cosh(\x)-0.5},{1.01*sinh(\x)});
\draw[dashed,blue] plot[domain=-2:2] ({1*cosh(\x)+0.75},{0.69*sinh(\x)});
\draw[dashed,blue] plot[domain=-2:2] ({-1*cosh(\x)-0.75},{0.69*sinh(\x)});
\draw[Mapcolor](0,3)--(0,-3);
\draw[blue](-3,3)--(3,3)--(3,-3)--(-3,-3)--(-3,3);
\tkzLabelPoint[below](s1){$A$}
\tkzLabelPoint[below](s2){$B$}
\end{scope}
\end{tikzpicture}}
\begin{itemize}
\item Xét hai nguồn sóng đồng bộ có phương trình lần lượt là: \bth{u_1=u_2=A\cos \omega t}.
\item Điểm M trong vùng giao thoa có khoảng cách tới các nguồn là \bth{d_1,\ d_2}.
\item Phương trình sóng tại điểm M do hai nguồn truyền tới
\begin{align*}
u_{1M}=A\cos(\omega t-2\pi\dfrac{d_1}{\lambda});\quad u_{2M}=A\cos(\omega t-2\pi\dfrac{d_2}{\lambda}).
\end{align*}
\item Phương trình giao thoa sóng tại M
\begin{align*}
u_M=u_{1M}+u_{2M}=2A\cos(\pi\dfrac{d_1-d_2}{\lambda})\cos(\omega t-\pi\dfrac{d_1+d_2}{\lambda}).
\end{align*}
\item Biên độ sóng tại M: \bth{\boder{A_M=2A\abs{\cos(\pi\dfrac{d_1-d_2}{\lambda})}}}.
\end{itemize}
\subsubsection{Vị trí cực đại và vị trí cực tiểu}
\immini[thm]{\begin{itemize}
\item Biên độ dao động tổng hợp đạt cực đại khi: 
\begin{align*}
\boder{d_1-d_2=k\lambda \quad \ct{với } k =0,\pm 1,\pm 2,\ldots} \end{align*}
\textit{Điểm M cực đại nếu hiệu đường đi tới hai nguồn bằng số nguyên lần bước sóng}
\item Biên độ dao động tổng hợp đạt cực tiểu khi:
\begin{align*}
\boder{d_1-d_2=\pbrace{k+\dfrac{1}{2}}\lambda \quad \ct{với } k =0,\pm 1,\pm 2,\ldots}
\end{align*}
\textit{Điểm M cực tiểu nếu hiệu đường đi tới hai nguồn bằng số bán nguyên lần bước sóng}
\end{itemize}
\newghichu{Trong các công thức ta có thể thay $d_1-d_2$ thành $d_2-d_1$, khi đó hình ảnh giao thoa sẽ lật \textit{"trái sang phải"}. Mọi kết quả đều không đổi}}
{\begin{tikzpicture}[draw=Mapcolor, very thick,]
\node[above] at (0,3){$0 \lambda$};
\node[above] at (1,3){$1 \lambda$};
\node[above] at (2,3){$2 \lambda$};
\node[above] at (-1,3){$-1 \lambda$};
\node[above] at (-2,3){$-2 \lambda$};
\node[below] at (0.5,-3){$0,5 \lambda$};
\node[below] at (1.7,-3){$1,5 \lambda$};
\node[below] at (-0.5,-3){$-0,5 \lambda$};
\node[below] at (-1.5,-3){$-1,5 \lambda$};
\draw[decorate,decoration={brace,mirror,amplitude=3pt},xshift=0pt,yshift=0pt]
(0.5,0) -- (1,0) node [black,midway,yshift=-15pt] 
{\footnotesize $\dfrac{\lambda}{2}$};
\begin{scope}[shift={(0,0)}]
\tkzInit[ymin=-3,ymax=3,xmin=-3,xmax=3]\tkzClip
\tkzDefPoints{-2/0/s1, 2/0/s2}
\tkzDrawPoints[size=3,fill=Mapcolor](s1,s2)
\draw[dashed,blue] plot[domain=-2:2] ({1*cosh(\x)-0.75},{4.5*sinh(\x)});
\draw[dashed,blue] plot[domain=-2:2] ({-1*cosh(\x)+0.75},{4.5*sinh(\x)});
\draw[Mapcolor] plot[domain=-2:2] ({1*cosh(\x)-0.5},{3*sinh(\x)});
\draw[Mapcolor] plot[domain=-2:2] ({-1*cosh(\x)+0.5},{3*sinh(\x)});
\draw[dashed,blue] plot[domain=-2:2] ({1*cosh(\x)-0.25},{2.2*sinh(\x)});
\draw[dashed,blue] plot[domain=-2:2] ({-1*cosh(\x)+0.25},{2.2*sinh(\x)});
\draw[Mapcolor] plot[domain=-2:2] ({1*cosh(\x)},{1.75*sinh(\x)});
\draw[Mapcolor] plot[domain=-2:2] ({-1*cosh(\x)},{1.75*sinh(\x)});
\draw[dashed,blue] plot[domain=-2:2] ({1*cosh(\x)+0.25},{1.35*sinh(\x)});
\draw[dashed,blue] plot[domain=-2:2] ({-1*cosh(\x)-0.25},{1.35*sinh(\x)});
\draw[Mapcolor] plot[domain=-2:2] ({1*cosh(\x)+0.5},{1.01*sinh(\x)});
\draw[Mapcolor] plot[domain=-2:2] ({-1*cosh(\x)-0.5},{1.01*sinh(\x)});
\draw[dashed,blue] plot[domain=-2:2] ({1*cosh(\x)+0.75},{0.69*sinh(\x)});
\draw[dashed,blue] plot[domain=-2:2] ({-1*cosh(\x)-0.75},{0.69*sinh(\x)});
\draw[Mapcolor](0,3)--(0,-3);
\draw[blue](s1)--(s2);
\tkzLabelPoint[left](s1){$A$}
\tkzLabelPoint[right](s2){$B$}
\end{scope}
\end{tikzpicture}}
